% Môn: Vật lý
% Lớp: 10
% Chương: Chương I. Mở đầu
% Bài: L10C1 Bài 2. Các quy tắc an toàn trong phòng thực hành Vật lí
% Số câu: 185
% App đọc file này trực tiếp — sửa rồi Commit trên GitHub, bấm Đồng bộ GitHub .tex

% L10C1 Bài 2. Các quy tắc an toàn trong phòng thực hành Vật lí
\dangbt{Nhận biết nội quy và quy tắc an toàn}
\begin{ex}
Khi tiến hành thí nghiệm, cần phải
\choice
{thảo luận nhóm để thống nhất quy tắc riêng của nhóm, có thể bỏ qua quy tắc an toàn của phòng thí nghiệm.}
{tiến hành thí nghiệm với thời gian ngắn nhất, không cần tuân thủ các quy tắc của phòng thí nghiệm.}
{tự đề xuất các quy tắc thí nghiệm để có thể tiến hành thí nghiệm nhanh nhất.}
{\True tuân theo các quy tắc an toàn của phòng thí nghiệm, hướng dẫn của giáo viên.}
\loigiai{Khi tiến hành thí nghiệm, cần phải tuân theo các quy tắc an toàn của phòng thí nghiệm, hướng dẫn của giáo viên.}
\end{ex}

\dangbt{An toàn khi sử dụng dụng cụ thủy tinh, vật sắc nhọn}
\begin{ex}
Chọn câu sai. Khi sử dụng các thiết bị quang học ta cần chú ý đến những điều gì?
\choice
{Sử dụng các thiết bị nhẹ nhàng.}
{Lau chùi cẩn thận thiết bị trước khi sử dụng.}
{Bảo quản nơi khô thoáng, tránh ẩm mốc.}
{\True Chần qua nước sôi để khử trùng.}
\loigiai{Các thiết bị quang học thường có các bộ phận nhạy cảm như thấu kính, gương. Việc chần qua nước sôi có thể làm hỏng các bộ phận này, đặc biệt là lớp phủ chống phản xạ hoặc các chi tiết nhựa, keo dán.}
\end{ex}

\dangbt{Nhận biết nội quy và quy tắc an toàn}
\begin{ex}
Chọn đáp án sai khi nói về những quy tắc an toàn khi làm việc với phóng xạ?
\choice
{Sử dụng phương tiện phòng hộ cá nhân như quần áo phòng hộ, mũ, găng tay, áo chì.}
{Tẩy xạ khi bị nhiễm bẩn phóng xạ theo quy định.}
{Kiểm tra sức khỏe định kì.}
{\True Ăn uống, trang điểm trong phòng làm việc có chứa chất phóng xạ.}
\loigiai{Ăn uống, trang điểm trong phòng làm việc có chứa chất phóng xạ là hành động sai, có thể gây nhiễm xạ.}
\end{ex}

\dangbt{Giải nghĩa các kí hiệu kĩ thuật và biển cảnh báo an toàn trong phòng thực hành}
\begin{ex}
Kí hiệu DC hoặc dấu “-” mang ý nghĩa là
\choice
{\True dòng điện 1 chiều.}
{dòng điện xoay chiều.}
{cực dương.}
{cực âm.}
\loigiai{Kí hiệu DC hoặc dấu “-” mang ý nghĩa là dòng điện 1 chiều (Direct Current).}
\end{ex}

\dangbt{Nhận biết nội quy và quy tắc an toàn}
\begin{ex}
Phát biểu nào sau đây là đúng khi nói về quy tắc an toàn trong phòng thí nghiệm?
\choice
{Tắt công tắc nguồn thiết bị điện sau khi cắm hoặc tháo thiết bị điện.}
{Tuyệt đối không tiếp xúc với các vật và các thiết bị thí nghiệm có nhiệt độ cao ngay khi có dụng cụ bảo hộ.}
{Được phép tiến hành thí nghiệm khi đã mang đồ bảo hộ.}
{\True Phải vệ sinh, sắp xếp gọn gàng, các thiết bị và dụng cụ thí nghiệm, bỏ chất thải thí nghiệm vào đúng nơi quy định sau khi tiến hành thí nghiệm.}
\loigiai{Phát biểu đúng là: Phải vệ sinh, sắp xếp gọn gàng, các thiết bị và dụng cụ thí nghiệm, bỏ chất thải thí nghiệm vào đúng nơi quy định sau khi tiến hành thí nghiệm.}
\end{ex}

\dangbt{Nhận biết nội quy và quy tắc an toàn}
\begin{ex}
Phát biểu nào sau đây là sai khi nói về những quy tắc an toàn trong phòng thí nghiệm?
\choice
{Kiểm tra cẩn thận thiết bị, phương tiện, dụng cụ thí nghiệm trước khi sử dụng.}
{Chỉ tiến hành thí nghiệm khi được sự cho phép của giáo viên hướng dẫn thí nghiệm.}
{Đọc kĩ hướng dẫn sử dụng thiết bị và quan sát các chỉ dẫn, các kí hiệu trên các thiết bị thí nghiệm.}
{\True Tắt công tắc nguồn thiết bị điện sau khi cắm hoặc tháo thiết bị điện.}
\loigiai{Phát biểu sai là: Tắt công tắc nguồn thiết bị điện sau khi cắm hoặc tháo thiết bị điện. Đúng phải là tắt công tắc nguồn thiết bị điện TRƯỚC KHI cắm hoặc tháo thiết bị điện.}
\end{ex}

\dangbt{Giải nghĩa các kí hiệu kĩ thuật và biển cảnh báo an toàn trong phòng thực hành}
\begin{ex}
Kí hiệu AC hoặc dấu "~" mang ý nghĩa là
\choice
{cực dương.}
{dòng điện 1 chiều.}
{cực âm.}
{\True dòng điện xoay chiều.}
\loigiai{Kí hiệu AC hoặc dấu "~" mang ý nghĩa là dòng điện xoay chiều (Alternating Current).}
\end{ex}

\dangbt{Nhận biết nội quy và quy tắc an toàn}
\begin{ex}
Quy tắc an toàn trong phòng thí nghiệm là phải
\choice
{tắt công tắc nguồn thiết bị điện sau khi cắm hoặc tháo thiết bị điện.}
{\True phải vệ sinh, sắp xếp gọn gàng, các thiết bị và dụng cụ thí nghiệm, bỏ chất thải thí nghiệm vào đúng nơi quy định sau khi tiến hành thí nghiệm.}
{được phép tiến hành thí nghiệm khi đã mang đồ bảo hộ.}
{tuyệt đối không tiếp xúc với các vật và các thiết bị thí nghiệm có nhiệt độ cao ngay khi có dụng cụ bảo hộ.}
\loigiai{Quy tắc an toàn trong phòng thí nghiệm là phải vệ sinh, sắp xếp gọn gàng, các thiết bị và dụng cụ thí nghiệm, bỏ chất thải thí nghiệm vào đúng nơi quy định sau khi tiến hành thí nghiệm.}
\end{ex}

\dangbt{Nhận biết và giải thích biển báo, kí hiệu cảnh báo}
\begin{ex}
Chọn đáp án sai. Cần tuân thủ các biển báo an toàn trong phòng thực hành nhằm mục đích
\choice
{chống cháy, nổ.}
{\True tạo ra nhiều sản phẩm mang lại lợi nhuận}
{hạn chế các trường hợp nguy hiểm như đứt tay, ngộ độc,…}
{tránh được các tổn thất về tài sản nếu không làm theo hướng dẫn.}
\loigiai{Mục đích của việc tuân thủ các biển báo an toàn là để đảm bảo an toàn cho người và tài sản, không phải để tạo ra lợi nhuận.}
\end{ex}

\dangbt{Nhận biết nội quy và quy tắc an toàn}
\begin{ex}
Hành động nào sau đây không phù hợp với các quy tắc an toàn trong phòng thực hành?
\choice
{\True Nếm thử để phân biệt các loại hóa chất.}
{Thu dọn phòng thực hành, rửa sạch tay sau khi đã thực hành xong.}
{Chỉ tiến hành thí nghiệm khi có người hướng dẫn.}
{Mặc đồ bảo hộ, đeo kính, khẩu trang.}
\loigiai{Nếm thử hóa chất là hành động cực kỳ nguy hiểm và không được phép trong phòng thực hành.}
\end{ex}

\dangbt{Quy trình sơ cứu khẩn cấp khi xảy ra tai nạn điện giật}
\begin{ex}
Nếu thấy có người bị điện giật chúng ta không được
\choice
{chạy đi gọi người tới cứu chữa.}
{tách người bị giật ra khỏi nguồn điện bằng dụng cụ cách điện.}
{ngắt nguồn điện.}
{\True dùng tay để kéo người bị giật ra khỏi nguồn điện.}
\loigiai{Khi thấy người bị điện giật, tuyệt đối không được dùng tay không để kéo người bị giật ra khỏi nguồn điện vì có thể bị điện giật theo.}
\end{ex}

\dangbt{Phân tích mối nguy và đánh giá mức độ rủi ro}
\begin{ex}
Thao tác nào sử dụng thiết bị thí nghiệm có thể gây nguy hiểm trong phòng thực hành?
\choice
{Để chất dễ cháy gần thí nghiệm mạch điện.}
{Không đeo găng tay cao su chịu nhiệt khi làm thí nghiệm với nhiệt độ cao.}
{Sử dụng dây điện đã bị sờn, cũ.}
{\True Tất cả các ý trên.}
\loigiai{Tất cả các thao tác trên đều có thể gây nguy hiểm trong phòng thực hành.}
\end{ex}

\dangbt{Nhận biết nội quy và quy tắc an toàn}
\begin{ex}
Quy tắc nào sau đây không phải là quy tắc an toàn trong thực hành?
\choice
{Đọc kỹ hướng dẫn sử dụng thiết bị và quan sát các chỉ dẫn, các kí hiệu trên các thiết bị thí nghiệm}
{\True Khi vào phòng thí nghiệm là thực hiện luôn thí nghiệm}
{Tắt công tắc ngồn thiết bị điện trước khi tháo hoặc cắm thiết bị điện}
{Phải bố trí dây điện gọn gàn, không bị vướn thi qua lại}
\loigiai{Khi vào phòng thí nghiệm cần đọc kỹ hướng dẫn, chuẩn bị đầy đủ và chỉ tiến hành thí nghiệm khi có sự cho phép của giáo viên, không được thực hiện luôn.}
\end{ex}

\dangbt{Nhận biết nội quy và quy tắc an toàn}
\begin{ex}
Quy tắc nào sau đây không đảm bảo an toàn trong phòng thực hành?
\choice
{\True Tiếp xúc với nơi có cảnh báo nguy hiểm về điện.}
{Tắc công tắc nguồn thiết bị trước khi cắm điện và sau khi tháo điện.}
{Tuân thủ sự hướng dẫn của giáo viên hướng dẫn.}
{Đọc kĩ hướng dẫn sử dụng thiết bị.}
\loigiai{Tiếp xúc với nơi có cảnh báo nguy hiểm về điện là hành động không đảm bảo an toàn.}
\end{ex}

\dangbt{Nhận biết nội quy và quy tắc an toàn}
\begin{ex}
Quy tắc nào sau đây không phải là quy tắc an toàn trong phòng thực hành Vật lí?
\choice
{Tắt công tắc nguồn thiết bị điện trước khi cắm hoặc tháo thiết bị điện.}
{\True Tiếp xúc trực tiếp với các vật và các thiết bị thí nghiệm có nhiệt độ cao.}
{Chỉ tiến hành thí nghiệm khi được sự cho phép của giáo viên hướng dẫn thí nghiệm.}
{Kiểm tra cẩn thận thiết bị, phương tiện, dụng cụ thí nghiệm trước khi sử dụng.}
\loigiai{Tiếp xúc trực tiếp với các vật và các thiết bị thí nghiệm có nhiệt độ cao là hành động nguy hiểm, không phải quy tắc an toàn.}
\end{ex}

\dangbt{Phân tích mối nguy và đánh giá mức độ rủi ro}
\begin{ex}
Phát biểu nào sau đây là sai khi nói về nguy cơ mất an toàn trong sử dụng thiết bị thí nghiệm Vật lí?
\choice
{\True Nguy cơ gây tật cận thị ở mắt.}
{Nguy cơ gây nguy hiểm cho người sử dụng.}
{Nguy cơ hỏng thiết bị đo điện.}
{Nguy cơ cháy nổ trong phòng thực hành.}
\loigiai{Nguy cơ gây tật cận thị ở mắt không phải là nguy cơ mất an toàn trực tiếp khi sử dụng thiết bị thí nghiệm Vật lí, mà thường liên quan đến việc đọc sách, làm việc ở khoảng cách gần trong thời gian dài.}
\end{ex}

\dangbt{Nhận biết nội quy và quy tắc an toàn}
\begin{ex}
Quy tắc nào sau đây không phải là quy tắc an toàn trong phòng thực hành?
\choice
{Phải bố trí dây điện gọn gàng, không bị vướng khi qua lại, tuyệt đối tuân thủ hướng dẫn của giáo viên.}
{\True Khi vào phòng thí nghiệm cần khẩn trương thực hiện luôn thí nghiệm.}
{Tắt công tắc nguồn thiết bị điện trước khi cắm hoặc tháo thiết bị điện.}
{Đọc kĩ hướng dẫn sử dụng thiết bị và quan sát các chỉ dẫn, các kí hiệu trên các thiết bị thí nghiệm.}
\loigiai{Khi vào phòng thí nghiệm cần chuẩn bị kỹ lưỡng, đọc hướng dẫn và chỉ thực hiện thí nghiệm khi có sự cho phép, không được khẩn trương thực hiện luôn.}
\end{ex}

\dangbt{Nhận biết nội quy và quy tắc an toàn}
\begin{ex}
Quy tắc nào sau đây là một trong các quy tắc an toàn trong phòng thí nghiệm?
\choice
{\True Kiểm tra cẩn thận thiết bị, phương tiện, dụng cụ thí nghiệm trước khi sử dụng.}
{Kiểm tra cẩn thận thiết bị, phương tiện, dụng cụ thí nghiệm sau khi sử dụng.}
{Không nhất thiết kiểm tra thiết bị, phương tiện, dụng cụ thí nghiệm khi trước sử dụng.}
{Kiểm tra thiết bị, phương tiện, dụng cụ thí nghiệm trước khi sử dụng.}
\loigiai{Kiểm tra cẩn thận thiết bị, phương tiện, dụng cụ thí nghiệm trước khi sử dụng là một quy tắc an toàn quan trọng để đảm bảo thiết bị hoạt động tốt và an toàn.}
\end{ex}

\dangbt{Nhận biết nội quy và quy tắc an toàn}
\begin{ex}
Hành động nào không tuân thủ quy tắc an toàn trong phòng thực hành?
\choice
{\True Dùng tay không để làm thí nghiệm.}
{Trước khi làm thí nghiệm với bình thủy tinh, cần kiểm tra bình có bị nứt vỡ hay không.}
{Bố trí dây điện gọn gàng.}
{Trước khi cắm, tháo thiết bị điện, sẽ tắt công tắc nguồn.}
\loigiai{Dùng tay không để làm thí nghiệm, đặc biệt là với hóa chất hoặc thiết bị nóng/điện, là hành động không tuân thủ quy tắc an toàn.}
\end{ex}

\dangbt{An toàn khi sử dụng thiết bị điện}
\begin{ex}
Điều nào sau đây có thể gây mất an toàn khi sử dụng thiết bị thí nghiệm?
\choice
{Cầm vào phần vỏ nhựa của đầu phích cắm để cắm vào ổ điện.}
{\True Tay ướt cầm vào đầu dây không có phích cắm để cắm vào ổ điện.}
{Đeo khẩu trang, găng tay khi thực hành thí nghiệm với hóa chất.}
{Sắp xếp thiết bị vào đúng vị trí sau khi sử dụng.}
\loigiai{Tay ướt cầm vào đầu dây không có phích cắm để cắm vào ổ điện có thể gây điện giật do nước là chất dẫn điện.}
\end{ex}

\dangbt{Nhận biết nội quy và quy tắc an toàn}
\begin{ex}
Phát biểu sai khi nói về những quy tắc an toàn khi làm việc với phóng xạ là
\choice
{giảm thời gian tiếp xúc với nguồn phóng xạ.}
{tăng khoảng cách từ ta đến nguồn phóng xạ.}
{đảm bảo che chắn những cơ quan trọng yếu của cơ thể.}
{\True không cần mang áo phòng hộ và không cần đeo mặt nạ.}
\loigiai{Khi làm việc với phóng xạ, việc mang áo phòng hộ và đeo mặt nạ là rất cần thiết để bảo vệ cơ thể khỏi tác hại của phóng xạ.}
\end{ex}

\dangbt{Nhận biết nội quy và quy tắc an toàn}
\begin{ex}
Hành động nào không tuân thủ quy tắc an toàn trong phòng thực hành?
\choice
{Trước khi cắm hay tháo thiết bị điện cần tắt công tắc nguồn.}
{\True Dùng tay không để cầm ống nghiệm hoặc bình đong.}
{Trước khi làm thí nghiệm với bình thủy tinh, cần kiểm tra bình có bị nứt vỡ hay không.}
{Bố trí dây điện gọn gàng.}
\loigiai{Dùng tay không để cầm ống nghiệm hoặc bình đong, đặc biệt khi chứa hóa chất nóng hoặc nguy hiểm, là hành động không tuân thủ quy tắc an toàn.}
\end{ex}

\dangbt{Phân tích mối nguy và đánh giá mức độ rủi ro}
\begin{ex}
Hoạt động nào sau đây có thể là nguyên nhân gây mất an toàn trong khi thực hiện các thí nghiệm thực hành?
\choice
{\True Dùng tay không cầm ống chứa hóa chất.}
{Sử dụng trang phục và trang thiết bị phù hợp với thí nghiệm.}
{Nhận biết các dụng cụ nguy hiểm trước khi tiến hành thí nghiệm.}
{Thực hiện đúng nội quy phòng thực hành.}
\loigiai{Dùng tay không cầm ống chứa hóa chất có thể gây bỏng, kích ứng da hoặc ngộ độc nếu hóa chất độc hại.}
\end{ex}

\dangbt{Nhận biết nội quy và quy tắc an toàn}
\begin{ex}
Trong các hoạt động dưới đây
1. Dùng tay ướt cắm điện vào nguồn điện.
2. Mang đồ ăn, thức uống, chạy nhảy, vui đùa khi vào phòng thí nghiệm.
3. Bỏ chất thải thí nghiệm vào đúng nơi quy định trong phòng thí nghiệm.
4. Rửa sạch da khi tiếp xúc với hóa chất.
5. Buộc tóc gọn gàng, tránh để tóc tiếp xúc với hóa chất và dụng cụ thí nghiệm.
6. Thực hiện thí nghiệm nhanh và mạnh.
Những hoạt động đảm bảo an toàn khi vào phòng thí nghiệm là
\choice
{\True 3, 4, 5.}
{1, 3, 5.}
{2, 4, 6.}
{1, 2, 7.}
\loigiai{Các hoạt động đảm bảo an toàn là: 3. Bỏ chất thải thí nghiệm vào đúng nơi quy định; 4. Rửa sạch da khi tiếp xúc với hóa chất; 5. Buộc tóc gọn gàng, tránh để tóc tiếp xúc với hóa chất và dụng cụ thí nghiệm.}
\end{ex}

\dangbt{Nhận biết nội quy và quy tắc an toàn}

\dangbt{Phân tích mối nguy và đánh giá mức độ rủi ro}

\dangbt{Nhận biết nội quy và quy tắc an toàn}
\begin{ex}
Trong các quy tắc dưới đây, quy tắc nào không phù hợp khi vào phòng thí nghiệm thực hành Vật lí?
\choice
{\True Chỉ tiến hành thí nghiệm khi có người khác cùng tham gia.}
{Giữ khoảng cách an toàn khi tiến hành thí nghiệm có nung nóng các vật hay các thí nghiệm có các vật bắn ra tia laser.}
{Kiểm tra cẩn thận thiết bị, phương tiện, dụng cụ thí nghiệm trước khi sử dụng.}
{Đọc kĩ hướng dẫn sử dụng và quan sát các chỉ dẫn, kí hiệu trên các thiết bị thí nghiệm.}
\loigiai{Quy tắc "Chỉ tiến hành thí nghiệm khi có người khác cùng tham gia" là không phù hợp. Quan trọng là phải có sự cho phép và hướng dẫn của giáo viên, không nhất thiết phải có người khác cùng tham gia.}
\end{ex}

\dangbt{Nhận biết nội quy và quy tắc an toàn}

\dangbt{Quy trình sơ cứu khẩn cấp khi xảy ra tai nạn điện giật}
\begin{ex}
Khi phát hiện người bị điện giật, ta phải làm gì đầu tiên?
\choice
{Gọi người đến sơ cứu.}
{Đưa người đó ra khỏi khu vực có điện.}
{Gọi cấp cứu.}
{\True Ngắt nguồn điện.}
\loigiai{Khi phát hiện người bị điện giật, việc đầu tiên cần làm là ngắt nguồn điện để đảm bảo an toàn cho cả người bị nạn và người cứu hộ.}
\end{ex}

\dangbt{Vận dụng quy tắc an toàn khi sử dụng các thiết bị đo điện, nhiệt và thủy tinh}
\begin{ex}
Để sử dụng an toàn thiết bị đo điện khi sử dụng ta cần
\choice
{không chọn đúng thang đo, nhầm lẫn thao tác.}
{chọn đúng thang đo, nhầm lẫn thao tác.}
{không chọn đúng thang đo, thực hiện đúng thao tác.}
{\True chọn đúng thang đo, thực hiện đúng thao tác.}
\loigiai{Để sử dụng an toàn thiết bị đo điện, cần chọn đúng thang đo và thực hiện đúng thao tác để tránh làm hỏng thiết bị và đảm bảo an toàn cho người sử dụng.}
\end{ex}

\dangbt{Nhận biết nội quy và quy tắc an toàn}
\begin{ex}
Quy tắc nào sau đây không phải là quy tắc an toàn trong phòng thực hành Vật lí?
\choice
{Kiểm tra cẩn thận thiết bị, phương tiện, dụng cụ thí nghiệm trước khi sử dụng.}
{\True Cần kiểm tra và tiếp xúc trực tiếp với các thiết bị có nhiệt độ cao.}
{Tắt công tắc nguồn thiết bị điện trước khi cắm hoặc tháo thiết bị điện.}
{Chỉ tiến hành thí nghiệm khi được sự cho phép của giáo viên hướng dẫn thí nghiệm.}
\loigiai{Tiếp xúc trực tiếp với các thiết bị có nhiệt độ cao là hành động nguy hiểm, không phải quy tắc an toàn.}
\end{ex}

\dangbt{Quy trình sơ cứu khẩn cấp khi xảy ra tai nạn điện giật}
\begin{ex}
Khi phát hiện Tuấn bị điện giật, ta phải làm gì đầu tiên?
\choice
{Gọi người đến sơ cứu.}
{Đưa người bị điện giật ra khỏi khu vực có điện.}
{\True Ngắt nguồn điện.}
{Gọi cấp cứu.}
\loigiai{Khi phát hiện người bị điện giật, việc đầu tiên cần làm là ngắt nguồn điện để đảm bảo an toàn cho cả người bị nạn và người cứu hộ.}
\end{ex}

\dangbt{Nhận biết nội quy và quy tắc an toàn}

\dangbt{Nhận biết nội quy và quy tắc an toàn}
\begin{ex}
Phát biểu nào sau đây là sai khi nói về những quy tắc an toàn khi làm việc với phóng xạ?
\choice
{Giảm thời gian tiếp xúc với nguồn phóng xạ.}
{Tăng khoảng cách từ ta đến nguồn phóng xạ.}
{Đảm bảo che chắn những cơ quan trọng yếu của cơ thể.}
{\True Mang áo phòng hộ và không cần đeo mặt nạ.}
\loigiai{Khi làm việc với phóng xạ, việc mang áo phòng hộ và đeo mặt nạ là rất cần thiết để bảo vệ cơ thể khỏi tác hại của phóng xạ.}
\end{ex}

\dangbt{Nhận biết nội quy và quy tắc an toàn}
\begin{ex}
Nội dung nào dưới đây không phải là quy tắc an toàn trong phòng thực hành?
\choice
{Đọc kĩ hướng dẫn sử dụng thiết bị.}
{Tắt công tắc nguồn thiết bị điện trước khi cắm hoặc tháo thiết bị.}
{Phải bố trí dây điện gọn gàng không bị vướng khi qua lại.}
{\True Được phép tiếp xúc trực tiếp với các vật, các thiết bị thí nghiệm có nhiệt độ cao.}
\loigiai{Tiếp xúc trực tiếp với các vật, các thiết bị thí nghiệm có nhiệt độ cao là hành động nguy hiểm, không phải quy tắc an toàn.}
\end{ex}

\dangbt{An toàn khi sử dụng thiết bị điện}
\begin{ex}
Trong các hoạt động dưới đây, những hoạt động nào tuân thủ nguyên tắc an toàn khi sử dụng điện?
(1) Chạm tay trực tiếp vào ổ điện, dây điện trần hoặc dây dẫn điện bị hở.
(2) Thường xuyên kiểm tra tình trạng hệ thông đường điện và các đồ dùng điện.
(3) Kiểm tra mạch có điện bằng bút thử điện.
(4) Sửa chữa điện khi chưa ngắt nguồn điện.
(5) Đến gần nhưng không tiếp xúc với các máy biến thế và lưới điện cao áp.
(6) Bọc kĩ các dây dẫn điện bằng vật liệu các điện.
\choice
{(1) – (2) – (3) – (6).}
{\True (2) – (3) – (6).}
{(2) – (3) – (4) – (5).}
{(1) – (3) – (6).}
\loigiai{Các hoạt động tuân thủ nguyên tắc an toàn khi sử dụng điện là: (2) Thường xuyên kiểm tra tình trạng hệ thống đường điện và các đồ dùng điện; (3) Kiểm tra mạch có điện bằng bút thử điện; (6) Bọc kĩ các dây dẫn điện bằng vật liệu cách điện.}
\end{ex}

\dangbt{Nhận biết nội quy và quy tắc an toàn}
\begin{ex}
Khi nói về những quy tắc an toàn khi làm việc với phóng xạ, phát biểu nào sau đây là sai?
\choice
{Giảm thời gian tiếp xúc với nguồn phóng xạ.}
{Tăng khoảng cách từ ta đến nguồn phóng xạ.}
{Đảm bảo che chắn những cơ quan trọng yếu của cơ thể.}
{\True Mang áo phòng hộ và không cần đeo mặt nạ.}
\loigiai{Khi làm việc với phóng xạ, việc mang áo phòng hộ và đeo mặt nạ là rất cần thiết để bảo vệ cơ thể khỏi tác hại của phóng xạ.}
\end{ex}

\dangbt{Nhận biết và giải thích biển báo, kí hiệu cảnh báo}

\dangbt{Nhận biết nội quy và quy tắc an toàn}

\dangbt{Nhận biết nội quy và quy tắc an toàn}
\begin{ex}
Phát biểu nào sau đây là sai khi nói về những quy tắc an toàn trong phòng thí nghiệm?
\choice
{\True Tiếp xúc với các vật và các thiết bị thí nghiệm có nhiệt độ cao ngay khi có dụng cụ bảo hộ.}
{Tắt công tắc nguồn thiết bị điện trước khi cắm hoặc tháo thiết bị điện.}
{Chỉ cắm phích cắm của thiết bị điện vào ổ cắm khi hiệu điện thế của nguồn điện tương ứng với hiệu điện thế định mức của dụng cụ.}
{Phải bố trí dây điện gọn gàng, không bị vướng khi qua lại.}
\loigiai{Ngay cả khi có dụng cụ bảo hộ, việc tiếp xúc trực tiếp với các vật và thiết bị có nhiệt độ cao vẫn cần cẩn trọng và tuân thủ quy trình, không phải là "ngay khi có dụng cụ bảo hộ" là được phép tiếp xúc một cách vô tư.}
\end{ex}

\dangbt{Nhận biết nội quy và quy tắc an toàn}
\begin{ex}
Phát biểu nào sau đây là sai khi nói về những quy tắc an toàn trong phòng thí nghiệm?
\choice
{Không tiếp xúc trực tiếp với các vật và các thiết bị thí nghiệm có nhiệt độ cao khi không có dụng cụ bảo hộ.}
{Không để nước cũng như các dung dịch dẫn điện, dung dịch dễ cháy gần thiết bị điện.}
{\True Được phép tiến hành thí nghiệm khi đã mang đồ bảo hộ.}
{Giữ khoảng cách an toàn khi tiến hành thí nghiệm nung nóng các vật, thí nghiệm có các vật bắn ra, tia laser.}
\loigiai{Việc được phép tiến hành thí nghiệm không chỉ phụ thuộc vào việc mang đồ bảo hộ mà còn cần sự cho phép của giáo viên và tuân thủ các quy trình khác.}
\end{ex}

\dangbt{Nhận biết nội quy và quy tắc an toàn}
\begin{ex}
Phát biểu nào sau đây là đúng khi nói về những quy tắc an toàn trong phòng thí nghiệm?
\choice
{Tắt công tắc nguồn thiết bị điện sau khi cắm hoặc tháo thiết bị điện.}
{Tuyệt đối không tiếp xúc với các vật và các thiết bị thí nghiệm có nhiệt độ cao ngay khi có dụng cụ bảo hộ.}
{Được phép tiến hành thí nghiệm khi đã mang đồ bảo hộ.}
{\True Phải vệ sinh, sắp xếp gọn gàng, các thiết bị và dụng cụ thí nghiệm, bỏ chất thải thí nghiệm vào đúng nơi quy định sau khi tiến hành thí nghiệm.}
\loigiai{Phát biểu đúng là: Phải vệ sinh, sắp xếp gọn gàng, các thiết bị và dụng cụ thí nghiệm, bỏ chất thải thí nghiệm vào đúng nơi quy định sau khi tiến hành thí nghiệm.}
\end{ex}

\dangbt{Xử lí sự cố và sơ cứu trong phòng thực hành}
\begin{ex}
Khi gặp sự cố mất an toàn trong phòng thực hành, học sinh cần
\choice
{\True báo cáo ngay với giáo viên trong phòng thực hành.}
{tự xử lí và không báo với giáo viên.}
{nhờ bạn xử lí sử cố.}
{tiếp tục làm thí nghiệm.}
\loigiai{Khi gặp sự cố mất an toàn, học sinh cần báo cáo ngay với giáo viên để được hướng dẫn xử lý kịp thời và đúng cách.}
\end{ex}

\dangbt{Xử lí sự cố và sơ cứu trong phòng thực hành}
\begin{ex}
Trong bài thực hành có sử dụng mạch điện nhưng khi lắp ráp xong mạch điện, báo cáo giáo viên phụ trách rồi cắm vào nguồn điện nhưng mạch không vào điện thì học sinh cần
\choice
{kiểm tra lại mạch điện.}
{\True ngắt mạch điện ra khỏi nguồn sau đó kiểm tra mạch điện và nguồn điện.}
{kiểm tra nguồn điện.}
{ngắt mạch điện ra khỏi nguồn.}
\loigiai{Khi mạch điện không vào điện, việc đầu tiên là ngắt mạch điện ra khỏi nguồn để đảm bảo an toàn, sau đó mới tiến hành kiểm tra mạch điện và nguồn điện.}
\end{ex}

\dangbt{An toàn với nguồn nhiệt và nguy cơ cháy nổ}
\begin{ex}
Khi phòng thực hành xuất hiện cháy thì ta cần phải
\choice
{\True ngắt điện, di chuyển các chất dễ cháy ra ngoài và chống cháy lan, cứu người, cứu tài sản, dập tắt đám cháy.}
{chạy ra khỏi phòng, đi tìm thêm người đến dập đám cháy.}
{ngắt nguồn điện, dùng nước dập đám cháy.}
{dùng nước dập đám cháy.}
\loigiai{Khi phòng thực hành xuất hiện cháy, cần thực hiện các bước theo thứ tự ưu tiên: ngắt điện, di chuyển chất dễ cháy, chống cháy lan, cứu người, cứu tài sản và dập tắt đám cháy bằng phương tiện phù hợp.}
\end{ex}

\dangbt{Nhận biết nội quy và quy tắc an toàn}
\begin{ex}
Hoạt động nào sau đây không được làm sau khi kết thúc giờ thí nghiệm?
\choice
{vệ sinh sạch sẽ phòng thí nghiệm.}
{sắp xếp gọn gàng các thiết bị và dụng cụ thí nghiệm.}
{bỏ chất thải thí nghiệm vào nới quy định.}
{\True để các thiết bị nối với nguồn điện giúp duy trì năng lượng.}
\loigiai{Sau khi kết thúc thí nghiệm, cần ngắt kết nối các thiết bị khỏi nguồn điện để đảm bảo an toàn và tiết kiệm năng lượng.}
\end{ex}

\dangbt{Nhận biết nội quy và quy tắc an toàn}
\begin{ex}
Trong các hoạt động dưới đây, hoạt động nào tuân thủ theo nguyên tắc an toàn khi làm việc với các nguồn phóng xạ
\choice
{\True sử dụng phương tiện phòng hộ cá nhân như quần áo phòng hộ, găng tay, mũ, áo chì.}
{ăn uống, trang điểm trong phòng nơi có chất phóng xạ.}
{đổ rác thải phóng xạ ra khu vực rác thải sinh hoạt.}
{tiếp xúc trục tiếp với chất phóng xạ.}
\loigiai{Sử dụng phương tiện phòng hộ cá nhân là nguyên tắc quan trọng để đảm bảo an toàn khi làm việc với phóng xạ.}
\end{ex}

\dangbt{Nhận biết nội quy và quy tắc an toàn}
\begin{ex}
Hoạt động nào sau đây không thực hiện đúng quy tắc an toàn trong phòng thực hành?
\choice
{Đeo găng tay khi làm thí nghiệm.}
{Không ăn uống, đùa nghịch trong phòng thí nghiệm.}
{\True Để hóa chất không đúng nơi quy định sau khi làm xong thí nghiệm.}
{Làm thí nghiệm theo sự hướng dẫn của giáo viên.}
\loigiai{Để hóa chất không đúng nơi quy định sau khi làm xong thí nghiệm là hành động không đúng quy tắc an toàn, có thể gây nguy hiểm hoặc nhầm lẫn.}
\end{ex}

\dangbt{Nhận biết nội quy và quy tắc an toàn}
\begin{ex}
Khi nói về những quy tắc an toàn trong phòng thí nghiệm, phát biểu đúng là
\choice
{tắt công tắc nguồn thiết bị điện sau khi cắm hoặc tháo thiết bị điện.}
{\True phải vệ sinh, sắp xếp gọn gàng, các thiết bị và dụng cụ thí nghiệm, bỏ chất thải thí nghiệm vào đúng nơi quy định sau khi tiến hành thí nghiệm.}
{tuyệt đối không tiếp xúc với các vật và các thiết bị thí nghiệm có nhiệt độ cao ngay khi có dụng cụ bảo hộ.}
{được phép tiến hành thí nghiệm khi đã mang đồ bảo hộ.}
\loigiai{Phát biểu đúng là: Phải vệ sinh, sắp xếp gọn gàng, các thiết bị và dụng cụ thí nghiệm, bỏ chất thải thí nghiệm vào đúng nơi quy định sau khi tiến hành thí nghiệm.}
\end{ex}

\dangbt{Nhận biết nội quy và quy tắc an toàn}
\begin{ex}
Thao tác đúng khi sử dụng thiết bị thí nghiệm trong phòng thực hành là
\choice
{rút phích điện khi dây điện hở.}
{\True đeo găng tay cao su chịu nhiệt khi làm thí nghiệm ở nhiệt độ cao.}
{cắm phích điện vào ổ mà tay lại chạm vào phần kim loại của phích điện.}
{đun nước trên đèn cồn.}
\loigiai{Đeo găng tay cao su chịu nhiệt khi làm thí nghiệm ở nhiệt độ cao là thao tác đúng để bảo vệ bản thân.}
\end{ex}

\dangbt{Nhận biết nội quy và quy tắc an toàn}
\begin{ex}
Đâu là quy tắc an toàn trong phòng thực hành?
\choice
{\True Tắt công tắc nguồn thiết bị điện trước khi cầm hoặc tháo thiết bị điện.}
{Tiếp xúc trực tiếp với các vật và các thiết bị thí nghiệm có nhiệt độ cao khi không có dụng cụ hỗ trợ.}
{Tiến hành thí nghiệm khi chưa được sự cho phép của giáo viên hướng dẫn thí nghiệm.}
{Để nước cũng như các dung dịch dẫn điện, dung dịch dễ cháy gần thiết bị điện.}
\loigiai{Tắt công tắc nguồn thiết bị điện trước khi cầm hoặc tháo thiết bị điện là một quy tắc an toàn cơ bản để tránh nguy cơ điện giật.}
\end{ex}

\dangbt{Nhận biết nội quy và quy tắc an toàn}
\begin{ex}
Trong các hoạt động dưới đây, hoạt động nào dưới đây, hoạt động nào đảm bảo an toàn vào phòng thí nghiệm?
(1) Mặc áo blouse, mang bao tay, kính bảo hộ trước khi vào phòng thí nghiệm.
(2) Dùng tay ướt cắm điện vào nguồn điện.
(3) Buộc tóc gọn gàng, tránh để tóc tiếp xúc với hóa chất và dụng cụ thí nghiệm.
(4) Mang đồ ăn, thức uống vào phòng thí nghiệm.
(5) Thực hiện thí nghiệm nhanh và mạnh.
(6) Rửa sạch da khi tiếp xúc với hóa chất.
(7) Tự ý đem đồ thí nghiệm mang về nhà luyện tập.
(8) Bỏ chất thải thí nghiệm vào đúng nơi quy định.
\choice
{(2) – (4) – (7) – (8).}
{(3) – (4) – (5)– (6) – (8).}
{(2) – (3) – (4)– (7) – (8).}
{\True (1) – (3) – (6) – (8).}
\loigiai{Các hoạt động đảm bảo an toàn là: (1) Mặc áo blouse, mang bao tay, kính bảo hộ trước khi vào phòng thí nghiệm; (3) Buộc tóc gọn gàng, tránh để tóc tiếp xúc với hóa chất và dụng cụ thí nghiệm; (6) Rửa sạch da khi tiếp xúc với hóa chất; (8) Bỏ chất thải thí nghiệm vào đúng nơi quy định.}
\end{ex}

\dangbt{An toàn khi sử dụng thiết bị điện}
\begin{ex}
Phát biểu nào sau đây là sai khi nói về những quy tắc an toàn khi một số biện pháp an toàn khi sử dụng điện?
\choice
{\True Đảm bảo che chắn những cơ quan trọng yếu của cơ thể gây tổ thương não.}
{Lắp đặt vị trí cầu dao, cầu chì, công tắc, ổ điện đúng nơi quy định}
{Tránh xa nơi điện thế nguy hiểm.}
{Không dùng tay ướt hoặc nhiều mồ hôi khi sử dụng dây điện.}
\loigiai{Phát biểu "Đảm bảo che chắn những cơ quan trọng yếu của cơ thể gây tổ thương não" không liên quan trực tiếp đến các biện pháp an toàn khi sử dụng điện mà thường liên quan đến an toàn phóng xạ hoặc hóa chất.}
\end{ex}

\dangbt{Quy trình sơ cứu khẩn cấp khi xảy ra tai nạn điện giật}

\dangbt{Phân tích nguy cơ mất an toàn liên quan đến thiết bị nhiệt và thủy tinh}
\begin{ex}
Nhận định nào sau đây đúng?
\choice
{Dụng cụ thí nghiệm là bình thủy tinh cực kỳ bền nên không lo bị nút, vỡ.}
{\True Việc thực hiện sai thao tác có thể gây nguy hiểm cho người sử dụng.}
{Việc thực hiện sai thao tác cùng lắm là thiết bị sẽ không hoạt động, không gây nguy hiểm tới người sử dụng.}
{Dây điện bị sờn chỉ mất tính thẩm mỹ, ngoài ra không gây nguy hiểm cho người sử dụng.}
\loigiai{Việc thực hiện sai thao tác có thể gây nguy hiểm cho người sử dụng, ví dụ như gây bỏng, điện giật, cháy nổ, hoặc làm hỏng thiết bị.}
\end{ex}

\dangbt{Phân tích mối nguy và đánh giá mức độ rủi ro}
\begin{ex}
Hành động nào sau đây không gây nguy hiểm cho người làm thực hành thí nghiệm?
\choice
{Để các kẹp điện gần nhau.}
{Không đeo găng tay cao su khi thực hiện làm ths nghiệm với nhiệt độ cao.}
{Để cồn gần thí nghiệm mạch điện.}
{\True Khi thí nghiệm với ampe kế cần cắm dây đo vào chốt cắm phù hợp với chức năng đo.}
\loigiai{Khi thí nghiệm với ampe kế cần cắm dây đo vào chốt cắm phù hợp với chức năng đo là thao tác đúng, đảm bảo an toàn và độ chính xác của phép đo.}
\end{ex}

\dangbt{Nhận biết nội quy và quy tắc an toàn}
